\documentclass[a4paper]{article}
\usepackage[pdftex]{graphicx}
\usepackage[utf8]{inputenc}
\usepackage{enumerate}
\usepackage{siunitx}
\sisetup{locale=DE}
\usepackage{href-ul}
\hypersetup{
	colorlinks=true,
	linkcolor=blue,
	urlcolor=blue}
\usepackage{icomma}
\usepackage{amssymb}
\usepackage{geometry}
\geometry{a4paper, top=15mm, left=15mm, right=15mm, bottom=15mm,
	headsep=10mm, footskip=12mm}


\begin{document}
	%Title
	{\bf Mechanical waves }\\
	\begin{enumerate}[1.]
		%Task 1
		\item A water wave with straight wave fronts is generated in a wave trough by a harmoniously vibrating exciter. The wave movement has already covered all points on the water surface. The x-axis of a Cartesian coordinate system runs horizontally and perpendicular to the wave fronts, the y-axis of the coordinate system is directed vertically upwards. The exciter of the water wave is at the point $x_0=\SI{0}{\centi\meter}$ and is moving straight up through the zero position at time $t_0=0s$. The wave that propagates from the exciter in the positive x direction can be described by the following wave equation:
		$$y(t,x)=\SI{0.70}{\centi\meter} \cdot \sin{ \left[ 2 \pi \left( {t \over \SI{0.40}{\second }}-{x \over \SI{6,0}{\centi\meter}} \right) \right]}$$
		
		\begin{enumerate}[a)]
			\item Calculate the frequency $f$ of the excitation oscillation and the magnitude $c$ of the wave propagation speed from the data of the wave equation.
			\item Draw the instantaneous image of the wave at a scale of 1:1 for the time $t_1=\SI{0.30}{\second}$
			Range $ \SI{0}{\centi\meter} \le x \le \SI{12,0}{\centi\meter}$
			Let \item $P$ be a point on the water surface with the x-coordinate $x_P=\SI{10,0}{\centi\meter}$. For time $t_1=\SI{0.30}{\second}$, calculate the magnitude of the velocity of point $P$. Draw the speed $\vec{v}~ (\SI{0.30}{\second}; \SI{10.0}{\centi\meter})$ in the instantaneous image of subtask 1. b), by attaching the associated vector arrow to point $P$. Use the scale
			$\SI{5,0}{\centi\meter\second^{-1}} ~\hat{=} ~\SI{1}{\centi\meter}$
			\item Also draw the vector arrow for the velocity $\vec{v} ~(\SI{0.30}{\second}; \SI{1.5}{\centi\meter})$ using the one given in 1 . c) specified scale into the image of subtask 1. b).
		\end{enumerate}
		
		%Exercise 2
		\item The wave with straight wave fronts from 1. now hits a wall perpendicular to the direction of propagation of the wave with two narrow gaps $S_1$ and $S_2$. The distance between the gap centers is $d=\SI{20,0}{\centi\meter}$. The waves emanating from the two slits propagate at a speed of $c=\SI{15}{\centi\meter\second^{-1}}$. To simplify matters, it should be assumed that each of the two waves alone excites any point on the water surface behind the wall to an oscillation with the amplitude $A=\SI{0.70}{\centi\meter}$.
		
		\begin{enumerate}[a)]
			\item A point $Q$ on the water surface is at a distance of $\SI{16.0}{\centi\meter}$ from gap $S_1$ and at a distance of $\SI{28.0}{\centi\meter}$ from gap $S_2$. The waves emanating from the two columns have already captured the point $Q$. Describe the motion of point $Q$. Explain how the phenomenon described comes about.
			\item If you increase the frequency of the excitation oscillation from the value calculated under 1.a), the point $Q$ remains practically at rest at certain frequencies. Calculate the frequency $f_1$ at which this occurs the first time.
			\item The number of curves with maximum oscillation amplitude depends on the frequency of the exciter oscillation. Determine the frequency range for which exactly seven curves with maximum oscillation amplitude can be observed.
		\end{enumerate}
		
	\end{enumerate}
	\fbox{
		\begin{minipage}{0.5\textwidth}
			Additional worksheets are available at
			
			\href{http://www.okuyakl.com}{www.okuyakl.com}
		\end{minipage}
		\hfill
		\begin{minipage}{0.4\textwidth}
			\includegraphics[width=1.5 cm]{../../viecher/zwe03}
			\includegraphics[width=2 cm]{../../viecher/afanticon1}
			
	\end{minipage}}
	%End
\end{document}